% questions.tex —— 1985 年·上海卷·理工农医类 · 题目内容片段（无答案）
% 由 tools/generate_projects.py 依来源总账自 1985/1985_数学真题_上海理_试题.md 生成，勿手改（改源 MD 后重跑）


\begin{gktip}
\textbf{来源说明}：本题面依据固定源 PDF《1985 年上海市普通高等学校招生统一考试题目·数学（理工农医类）》第 1—2 页逐题转录。原卷页码为 4-36、4-37；该 PDF 不含答案页，下列答案和解析由本库独立推导，不声称代表官方发布版本。

\textbf{题面异文说明}：第 2 题第（5）小题固定源扫描字面印为“曲线 $x=x^3$”，但两个公开答案资料均将本卷第 2 题五小题列为（D），且只有将曲线核定为 $y=x^3$ 才能与四个选项的平移换元一致。规范源采用经交叉核定的 $y=x^3$；固定扫描的 $x=x^3$ 差异、来源 URL 与修订理由见问题总账 ISSUE-0763—ISSUE-0765，未修改上游 PDF。
\end{gktip}

\begin{question}
\gkti{1}（本题满分 30 分）

本题共有 10 个小题，每一个小题满分 3 分，只要求直接填写结果。

\Needspace{5\baselineskip}
（1） 不等式

\nopagebreak[4]
\[
|x+2|<1
\]

的解集是　　　　。

\Needspace{5\baselineskip}
（2） 函数

\nopagebreak[4]
\[
y=\sqrt[8]{x}-1
\]

的反函数是　　　　。

（3） 点 $(0,1)$ 到直线 $x+y=2$ 的距离是　　　　。

\Needspace{5\baselineskip}
（4） 函数

\nopagebreak[4]
\[
y=\sin\frac{\pi x}{2}
\]

的最小正周期是　　　　。

（5） 如果一个圆的圆心在点 $(2,4)$，并且经过点 $(0,3)$，那么这个圆的方程是　　　　。

（6） 若平面 $\alpha$ 及这个平面外的一条直线 $l$ 同时垂直于直线 $m$，则直线 $l$ 和平面 $\alpha$ 的位置关系是　　　　。

\Needspace{5\baselineskip}
（7） 函数

\nopagebreak[4]
\[
y=\sin^2x-\cos^2x
\]

的最小值是　　　　。

\Needspace{5\baselineskip}
（8） 方程

\nopagebreak[4]
\[
9^x-7\cdot3^x-18=0
\]

的解是　　　　。

（9） 若 $\operatorname{tg}\alpha=-2$，且 $\sin\alpha<0$，则 $\cos\alpha$ 的值是　　　　。

（10） 已知 $O$ 为直角坐标系的原点，点 $A$ 的坐标为 $(1,3)$，点 $B$ 的坐标为 $(0,3)$，若把 $\triangle OAB$ 绕 $y$ 轴旋转一周，则所得旋转体的体积是　　　　。

\end{question}
\begin{question}
\gkti{2}（本题满分 15 分）

本题共有 5 个小题，每一个小题都给出代号为 A、B、C、D 的四个结论，其中只有一个结论是正确的，把正确结论的代号写在题后的圆括号内，选对得 3 分，不选、选错或者选出的代号超过一个（不论是否都写在圆括号内），一律得 0 分。

（1） 设 $z$ 为复数，$z$ 是 $\overline z$ 的共轭复数，则 $z=\overline z$ 是 $z$ 为实数的

（A）充分但不必要的条件　　（B）必要但不充分的条件　　（C）充分而且必要的条件　　（D）既不充分也不必要的条件

（2） 下列各式中，正确的是

（A）$\arcsin\left(-\dfrac\pi3\right)=-\dfrac{\sqrt3}{2}$

（B）$\sin\left(\arcsin\dfrac\pi3\right)=\dfrac\pi3$

（C）$\arcsin\left(\sin\dfrac{5\pi}{4}\right)=-\dfrac\pi4$

（D）$\sin\left[\arccos\left(-\dfrac12\right)\right]=-\dfrac{\sqrt3}{2}$

\Needspace{5\baselineskip}
（3） 已知函数

\nopagebreak[4]
\[
f(x)=\lg(x^2-3x+2)
\]

\Needspace{5\baselineskip}
的定义域为 $F$，函数

\nopagebreak[4]
\[
g(x)=\lg(x-1)+\lg(x-2)
\]

的定义域为 $G$，那么

（A）$F\cap G=\varnothing$　　（B）$F=G$　　（C）$F\subset G$　　（D）$G\subset F$

（4） 若一个棱锥的底面是边数大于 3 的凸多边形，它的顶点到底面各边的距离都相等，则这个棱锥的底面多边形

（A）必为正多边形　　（B）必有内切圆　　（C）必有外接圆　　（D）必既有内切圆又有外接圆

（5） 若平移坐标轴，把坐标系 $xOy$ 的原点 $O$ 移到点 $O'$，$O'$ 在原坐标系中的坐标为 $(2,-1)$，则原坐标系中的曲线 $y=x^3$ 在新坐标系 $x'O'y'$ 中的方程是

（A）$y'+1=(x'-2)^3$　　（B）$y'+1=(x'+2)^3$

（C）$y'-1=(x'-2)^3$　　（D）$y'-1=(x'+2)^3$

\Needspace{0.4\textheight}
\end{question}
\begin{question}
\gkti{3}（本题满分 9 分）

本题共有 2 个小题，第（1）小题满分 5 分，第（2）小题满分 4 分。

\Needspace{5\baselineskip}
（1） 在直角坐标系内，方程

\nopagebreak[4]
\[
\left\{\begin{aligned}
x&=3\operatorname{ctg}\phi,\\
y&=2\operatorname{csc}\phi
\end{aligned}\right.
\]

（$\phi$ 是参数）表示什么曲线？画出它的图形。

\Needspace{5\baselineskip}
（2） 在极坐标系内，方程

\nopagebreak[4]
\[
\frac{2\sin\theta}{\rho}=1
\]

表示什么曲线？画出它的图形。
\nopagebreak[4]
\begin{center}% 1985 年上海卷理工农医类第 3 题：双曲线与圆
\begin{tikzpicture}[scale=0.72,>=stealth]
  \begin{scope}[shift={(-3.0,0)}]
    \draw[->] (-2.5,0)--(2.55,0) node[right] {$x$};
    \draw[->] (0,-2.7)--(0,2.7) node[above] {$y$};
    \draw[thick,domain=-2.2:2.2,samples=100]
      plot (\x,{2*sqrt(1+\x*\x/9)});
    \draw[thick,domain=-2.2:2.2,samples=100]
      plot (\x,{-2*sqrt(1+\x*\x/9)});
    \fill (0,2) circle (1.1pt);
    \fill (0,-2) circle (1.1pt);
    \node[below left] at (0,0) {$O$};
    \node[below] at (0,-3.0) {（1）};
  \end{scope}
  \begin{scope}[shift={(2.35,0)}]
    \draw[->] (-1.45,0)--(1.5,0) node[right] {$x$};
    \draw[->] (0,-0.45)--(0,2.55) node[above] {$y$};
    \draw[thick] (0,1) circle (1);
    \fill (0,1) circle (1.1pt);
    \node[below left] at (0,0) {$O$};
    \node[below] at (0,-0.75) {（2）};
  \end{scope}
\end{tikzpicture}
\end{center}
\nopagebreak[4]
\end{question}
\begin{question}
\gkti{4}（本题满分 16 分）

本题共有 3 个小题，第（1）、（2）小题满分各 5 分，第（3）小题满分 6 分。

\Needspace{5\baselineskip}
（1） 求

\nopagebreak[4]
\[
\left(x^3-\frac1{x^2}\right)^{15}
\]

的展开式中的常数项。

（2） 从六个数字 $1,2,3,4,5,6$ 中任取四个不同的数字，有多少种取法？由这六个数字可以组成多少个没有重复数字的四位偶数？

（3） 在一列球中，第 1 个球的半径为 1，第 2 个球的直径是第 1 个球的半径，以后每一个球的直径都是前一个球的半径，这样无限继续下去，求所有这些球的表面积的和。

\end{question}
\begin{question}
\gkti{5}（本题满分 6 分）

设 $\triangle ABC$ 的两个内角 $A\text{、}B$ 所对的边分别为 $a\text{、}b$。若复数 $z_1$ 的模为 $a$，幅角为 $B$；复数 $z_2$ 的模为 $b$，幅角为 $-A$。求证 $z_1+z_2$ 为实数。

\end{question}
\begin{question}
\gkti{6}（本题满分 8 分）

\Needspace{5\baselineskip}
解三角方程

\nopagebreak[4]
\[
\cos^2x+\cos^23x=1
\]

\end{question}
\begin{question}
\gkti{7}（本题满分 6 分）

已知函数 $f(x)$ 对其定义域内的任意两个实数 $a\text{、}b$，当 $a<b$ 时，都有 $f(a)<f(b)$。试用反证法证明：方程 $f(x)=0$ 至多有一个实数根。

\end{question}
\begin{question}
\gkti{8}（本题满分 8 分）

\Needspace{5\baselineskip}
对于一切大于 1 的自然数 $n$，证明：

\nopagebreak[4]
\[
\left(1+\frac13\right)\left(1+\frac15\right)\cdots\left(1+\frac1{2n-1}\right)>\frac{\sqrt{2n+1}}2
\]

\end{question}
\begin{question}
\gkti{9}（本题满分 10 分）

在三棱锥 $P-ABC$ 中，顶点 $P$ 到 $BC\text{、}AC\text{、}AB$ 的距离分别为 $h_1\text{、}h_2\text{、}h_3$，二面角 $P$- $BC$- $A=\alpha_1$，$P$- $AC$- $B=\alpha_2$，$P$- $AB$- $C=\alpha_3$。$\alpha_1\text{、}\alpha_2\text{、}\alpha_3$ 都是锐角，且依次成等差数列，$h_1\text{、}h_2\text{、}h_3$ 依次成等比数列。求证：$\alpha_1=\alpha_2=\alpha_3$。

\end{question}
\begin{question}
\gkti{10}（本题满分 12 分）

\Needspace{5\baselineskip}
已知：直线 $y=x+m$ 和曲线

\nopagebreak[4]
\[
x^2+2y^2+y^4-1=0
\]

交于 $A\text{、}B$ 两点，$P$ 是这条直线上的点，且 $|PA|\cdot|PB|=2$。求：当 $m$ 变化时，点 $P$ 的轨迹方程，并说明轨迹是什么图形。

\end{question}